% Lösungen Einheit 3 von 5 – Verflechtung (zusammenbau v0.1)
\einheitenkopf{Einheit 3 von 5 · Verflechtung}

\erg{13}{a) von R1 nach Z2 \quad b) $A \cdot (5 | 4) = (22 | 13)$: 22 kg Mehl, 13 kg Zucker \quad c) $A = ((2 | 5), (0 | 1), (4 | 0))$ – fehlende Pfeile sind Nullen \quad d) A: Zeilenzahl $= 4$, Spaltenzahl $= 3$; B: Zeilenzahl $= 3$, Spaltenzahl $= 2$; C: Zeilenzahl $= 4$, Spaltenzahl $= 2$ \quad e) $1 \cdot 1 + 2 \cdot 4 = 9$ – Pfad R1 → Z1 → E2 und Pfad R1 → Z2 → E2 \quad f) $x = 15$ Stück E1, $y = 20$ Stück E2 (aus $2x + y = 50$ und $x + 4y = 95$) \quad g) 30 Stück: R1 reicht für 40, R2 für 30, R3 für 45 – R2 ist der knappste Rohstoff \quad h) $z = 10$: Z1 kostet 10 GE, Z2 30 GE je ME. $B \cdot (20 | 10) = (50 | 30)$ sind die ME der Zwischenprodukte für den Auftrag, $(z | 3z)$ die Fertigungskosten je ME (Z2 dreimal so teuer wie Z1), 1400 GE die Fertigungskosten des Auftrags; aus $140z = 1400$}
\erg{14}{a) Fehler: nur mit der ersten Stufe gerechnet, A gilt je Zwischenprodukt. Richtig: $C = A \cdot B = ((1 | 5), (2 | 4))$, Bedarf $C \cdot (10 | 20) = (110 | 100)$ \quad b) Aus $r = A \cdot z$ und $z = B \cdot e$ folgt $r = A \cdot (B \cdot e) = (A \cdot B) \cdot e$; B muss zuerst auf die Endproduktmengen wirken, $B \cdot A$ passt nicht zur Folge der Stufen \quad c) ja: Bedarf $(900 | 600)$ l Apfel- und Traubensaft, das Lager hat 1000 l und 800 l}
